\scanpage{117}
% archive foliation: 116
Comportement fonctoriel des groupes d'inertie.

1. Soient $X$ un schéma, $U$ un ouvert de $X$, $G$ un groupe profini, $U'$ un revêtement principal de $U$ de groupe $G^{\deleted{\unsure{'}}}$, $x$ un point \inserted{géom.} de $X$, $\overline X$ le localisé strict de $X$ en $x$, $\overline U$ l'image inverse de $U$ dans $\overline X$, $\overline U'=U'\times_U\overline U$. On suppose donné de plus un schéma $\overline X'$ sur $\overline X$, avec des opérations de $G$ sur $\overline X'$, de telle façon que le \deleted{espace} \inserted{$\overline U$-schéma} à opérateurs $\overline X'|\overline U$ coïncide avec $\overline U'$.

Pratiquement, on sera intéressé par le cas où $x\in X-U$ et où $x$ est adhérent à $U$, où $\overline X'$ est entier sur $\overline X$ et où $\underline{\mathrm O}_{\overline X}=\overline p(\underline{\mathrm O}_{\overline X'})^G$, $\overline p:\overline X'\to\overline X$ étant la projection canonique, - de sorte que $G$ opère transitiv\deleted{\unsure{e}}ement sur l'ensemble fibre $\overline X'_{\overline x}$. De plus, on \deleted{\unsure{se placera}} sera \inserted{le plus souvent,} sous des conditions où on peut construire canoniquement $\overline X'$ en termes de $\overline U'$; par exemple, si on suppose $X$ normal en $x$, donc $\overline X^{\deleted{\unsure{'}}}$ normal, et $\overline U$ non vide, on définira $\overline X'$ comme le normalisé de $\overline U$ relativement à $\overline U'$ sur $\overline U$ (i.e. comme limite projective des \unsure{normaalisés|normalisés} dans les $\overline{U'(i)}$ sur $\overline U$, $U'$ étant écrit comme limite projective de schémas principaux $U'(i)$ sur les groupes quotients finis $G(i)$ de $G$).

Si maintenant \deleted{\unsure{$X$} est} \inserted{$\xi$} est un point de $\overline X'_{\overline x}$, on peut considérer le groupe de stabilité $G_\xi$ de $\xi$, formé des $g\in G$ tels que $g.\xi=\xi$; ce groupe prendra le nom de \uline{groupe d'inertie} \uline{relatif à} $\xi$. Dans le cas envisagé plus haut où les points de $\overline X'$ sur $\overline x$ sont tous conjugués sous $G$, on voit que pour $\xi$ variable\deleted{\unsure{s}}, les $G_\xi$ parcourent une classe de sous-groupes conjugués de $G$, qu'on appellera \inserted{(par abus de langage)} \uline{les sous-groupes d'inertie} \deleted{de} \uline{relatifs à} $x$. Noter qu'en toute rigueur cette classe de sous-groupes \deleted{d'inertie} n'est pas déterminée par la seule connaissance de $x\in X$, mais dépend du choix du prolongement de $\overline U'$ en $\overline X'$. Lorsque $X$ est normal en $x$ et que $x$ est spécialisation d'un point de $U$ (ce qui signifie que $x$ est adhérent à $U$ si $\operatorname{esp}(X)$ est localement noethérien), il sera sous-entendu, sauf mention du contraire, que le choix en question est fait par normalisation, comme expliqué plus haut.
